\documentclass[12pt,superscriptaddress]{revtex4}
%\documentclass[prl,superscriptaddress,twocolumn]{revtex4}
%\usepackage[lite,subscriptcorrection,slantedGreek,nofontinfo]{mtpro2}
%\usepackage{times}
%\usepackage{CJK}
%\usepackage{ctex}

\usepackage{amsfonts}
\usepackage{mathrsfs}
\usepackage{amsmath}
\usepackage{amssymb}
%\usepackarrge{mathptmx}
\usepackage{bm}

\usepackage{feynmp-auto}
\usepackage[normalem]{ulem}
\usepackage{extarrows}
\usepackage{slashed}
\usepackage{isodateo}
\usepackage{graphicx}
\usepackage{xcolor}
\usepackage[CJKbookmarks=true, bookmarksnumbered=true,bookmarksopen=true]{hyperref}
\hypersetup{colorlinks,%
	linkcolor=[rgb]{0,0.2,0.6}, %
	citecolor=[rgb]{0,0.2,0.6},
	urlcolor=[rgb]{0,0.2,0.6}}

\usepackage{indentfirst}

\graphicspath{{fig/}}


\linespread{1.1}
\renewcommand{\FR}[2]{\displaystyle\frac{\,{#1}\,}{#2}}
\newcommand{\fr}[2]{\mbox{$\frac{\,{#1}\,}{#2}$}}

\renewcommand{\thefootnote}{\arabic{footnote}}

\newcommand{\order}[1]{\mathcal{O}({#1})}



\newcommand{\del}[1]{{\blue\sout{#1}}}


\newcommand{\red}{\color{red}}
\newcommand{\blue}{\color{blue}}

\newcommand{\ob}[1]{\mkern 2mu \overline{\mkern -2mu #1 \mkern -2mu}\mkern 2mu}
\newcommand{\wt}[1]{\mkern 2mu \widetilde{\mkern -2mu #1 \mkern -2mu}\mkern 2mu}
\newcommand{\wh}[1]{\mkern 2mu \widehat{\mkern-2mu#1\mkern-2mu}\mkern 2mu}
\def\ma{\mathcal}





\begin{document}
	
	
	\title{Response to the Review Comments}
    \author{Jiaqi Chen}
    \author{Bo Feng}
	\author{Yi-Xiao Tao}
	
	
	\maketitle

\noindent Dear Editors and Reviewers:
	
We would like to thank you for your careful reading, helpful comments, and constructive suggestions, which has significantly improved the presentation of our manuscript. We have carefully considered all comments from the reviewers and revised our manuscript accordingly. The manuscript has also been double-checked, and the typos and grammar errors we found have been corrected. In the following section, we summarize our responses to each comment from the reviewers. We believe that our responses have well addressed all concerns from the reviewers. The main corrections in the paper and the responds to the reviewer's comments are as following:

\noindent\rule[0.25\baselineskip]{\textwidth}{1pt}

\begin{center}
\fcolorbox{black}{gray!10}
{
	\parbox{1\linewidth}
	{		
		 For the In-In formalism, the authors shall quote the classical references
 Keldysh:1964ud,Schwinger:1960qe,Chou:1984es, and B. DeWitt, in: Quantum Concepts in Space and Time, eds. R. Penrose and C.J. Isham, Calendron
 Press, Oxford, Calendron, 1986).

	}
}
\end{center}
\noindent\textcolor{red}{Response}:
Thank you for this suggestion. We have added these references in the second paragraph of our introduction.

\begin{center}
\fcolorbox{black}{gray!10}
{
	\parbox{1\linewidth}
	{		
		 Furthermore, they shall quote standard textbooks about hypergeometric
 functions and their extensions,such as: W.N. Bailey, Generalized Hypergeometric Series, (Cambridge University Press, Cambridge, 1935). L.J. Slater,
 Generalized hypergeometric functions, (Cambridge University Press, Cam
bridge, 1966). H. Exton, Multiple Hypergeometric Functions and Applica
tions, (Ellis Horwood, Chichester, 1976). H. Exton, Handbook of Hyperge
ometric Integrals, (Ellis Horwood, Chichester, 1978). H.M. Srivastava and
 P.W. Karlsson, Multiple Gaussian Hypergeometric Series, (Ellis Horwood,
 Chicester, 1985) because quite a series of representations and transformations they are mentioning are found there.

 As for the analytic continuations of these higher transcendental functions
 the background can be given by using Mellin-Barnes integral representations,
 which has to be mentioned in explicit form and by referring to Slater’s book
 1
 REPORT JHEP\_246P\_1124
e.g. More precisely, they should refer to Pochhammer’s contour integral,
 which might not be known to them. It is described in: Whittaker, E. T.;
 Watson, G. N. (1963), A Course of Modern Analysis, Cambridge University
 Press, ISBN 978-0-521-58807-2

	}
}
\end{center}
\noindent\textcolor{red}{Response}:
Thank you for pointing out this. We have added these references in section 3.1.4 under the equation (3.26).


\begin{center}
\fcolorbox{black}{gray!10}
{
	\parbox{1\linewidth}
	{		
  Difference and differential equations up to multivariate case of hypergeometric functions and their generalizations in QFT have also been discussed
 in Blumlein:2021hbq recently, giving different algorithms to deal with these
 structures.

 For (3.62) Bera shall not be quoted (he is a late user only), since F4 is
 known for much longer. I suggest Slater’s standard textbook here.

	}
}
\end{center}
\noindent\textcolor{red}{Response}:
Thank you for this useful suggestion. We have added this reference in footnote 1 of our introduction. We have also corrected the wrong reference.

\noindent\rule[0.25\baselineskip]{\textwidth}{1pt}




% We can also mention that the only difference between our results and former literature is that there is a term has redundant imaginary coefficient $i$ in their results.
We tried our best to improve the manuscript and made some changes in the manuscript. And we appreciate your warm work earnestly and hope the correction will meet with approval.

	
% \bibliographystyle{JHEP}
% \bibliography{Reply}
	
	
\end{document}
